\subsection{同志桑拿与性场所：文化、安全与法律}

俚语词条中的"SN（桑拿）"指向一个真实存在的场所类型：同志桑拿（gay sauna / bathhouse），以及会所、公园等偶遇性空间（cruising 空间）。理解它们的健康意义需要先理解其历史：在出柜可能带来失业与家庭破裂的年代，这类空间是男男社群几乎唯一的联结场所——它们承载的不只是性，还有社交、身份认同与互助网络。今天，场所的社交功能仍真实存在，健康议题也真实存在。

\textbf{健康要点}：
\begin{itemize}
\item \textbf{屏障规则在暗处更难维持}：暗房、迷宫等低照明场景中，"来不及戴套"是常见现实。可行做法是把安全套与水性润滑预先装在口袋里，进入暗区前完成佩戴与协商；对中途摘套的行为保持警觉——这是场所内高发的越界行为，也是安全套失败的主要原因之一；
\item \textbf{药物风险}：场所文化中"G 毒品""冰"等药物的使用与 chemsex 场景高度重叠（参见本书"性与药物"章节）——药物削弱的恰恰是拒绝能力与屏障坚持度；饮品不离视线、不单独进入封闭空间的规则同样适用；
\item \textbf{暴露后处理}：无保护肛交是 HIV 传播效率最高的行为之一，高危暴露后 72 小时内阻断（PEP）的时间窗在场所情境中尤其宝贵——平时就查好最近的阻断门诊在哪里，比事后搜索有用得多；
\item \textbf{定期检测}：性活跃的多伴侣人群建议每 3 个月检测一次（HIV、梅毒、淋病、衣原体）；许多城市的社群组织提供免费快速检测，场所周边常能获取相关信息。
\end{itemize}

\textbf{法律边界}：需要区分"场所内自愿私密的性行为"与"公共场所公开裸露"——后者可能触及治安管理处罚法的相关规定，场所管理方也负有合规责任。更重要的是防范侵害：场所内及周边的勒索、偷拍与暴力事件并不罕见。遇到侵权应保留证据并报警；性取向与性行为不是过错，勒索者才是违法者。当伴侣身份无法公开时，"被敲诈者往往不敢报警"正是勒索得以成立的机制——了解这一点，本身就是一种防护。
